% 第 12 讲 导出函子与显式投射消解
% 源：src/lecture12.txt（244 行）
\section*{第 12 讲\quad 导出函子与显式投射消解}

设 $X$ 与 $Y$ 是 $A$-模复形。回忆上一讲定义了 $\operatorname{Hom}_A(X,Y)$，以及
\[
\operatorname{Hom}^{\mathrm{der}}_A(X,Y):=\operatorname{Hom}_A(P,Y),
\]
其中 $P\xrightarrow{\mathrm{qis}}X$ 是投射复形，即 $X$ 的一个投射消解。我们还定义了 \textbf{Ext}-群
\[
\operatorname{Ext}^i_A(X,Y):=H^i\operatorname{Hom}^{\mathrm{der}}_A(X,Y).
\]
这里最重要的例子是 $A:=\Z[G]$，此时 $X^{hG}:=\operatorname{Hom}^{\mathrm{der}}_A(\Z,X)$ 就是 $X$ 的同伦不变量。这一构造具有如下基本性质：$\operatorname{Hom}^{\mathrm{der}}_A(X,-)$ 保持拟同构；且 $P\xrightarrow{\mathrm{qis}}X$ 在同伦意义下唯一，这样的同伦又在同伦意义下唯一，如此递推。

附带说明，我们其实可以更典范地定义导出函子 $\operatorname{Hom}^{\mathrm{der}}(X,-)$，而不必选取特定的投射消解：
\[
\operatorname{Hom}^{\mathrm{der}}(X,Y):=\varinjlim_{P\xrightarrow{\mathrm{qis}}X,\ P\ \text{投射}}\operatorname{Hom}(P,Y),
\]
其中 $P$ 按链复形映射 $P'\xrightarrow{\mathrm{qis}}P\xrightarrow{\mathrm{qis}}X$ 排序，这便强制消除了对 $P$ 的选取。

\textbf{断言 12.1.} 设 $f\colon X_1\to X_2$ 是链复形映射，$X_1$、$X_2$ 分别有投射消解 $P_1$、$P_2$。则存在映射 $\phi\colon P_1\to P_2$，使下图在同伦意义下交换：
\[
\begin{array}{ccc}
X_1 & \xrightarrow{\ f\ } & X_2\\
{\scriptstyle\mathrm{qis}}\big\uparrow & & \big\uparrow{\scriptstyle\mathrm{qis}}\\
P_1 & \xrightarrow{\ \phi\ } & P_2
\end{array}
\]
此外，这样的映射在同伦意义下唯一。

\textit{证明.} 由于导出函子保持拟同构，诱导出的映射复形之间的映射
\[
\operatorname{Hom}(P_1,P_2)\xrightarrow{\mathrm{qis}}\operatorname{Hom}(P_1,X_2)
\]
是拟同构。我们已给定一个映射，即复合 $P_1\to X_1\to X_2$；它被微分杀死（因为它是链复形映射），从而定义了 $H^0\operatorname{Hom}(P_1,X_2)$ 中的一个上同调类。于是 $H^0\operatorname{Hom}(P_1,P_2)$ 中存在某个上同调类，它是该映射经 $P_2$ 的提升，且良定义、在同伦意义下唯一。\hfill$\square$

由此可见，对每个链复形 $Y$，沿 $\phi$ 拉回便得到映射
\[
\operatorname{Hom}^{\mathrm{der}}_A(X_2,Y)\to\operatorname{Hom}^{\mathrm{der}}_A(X_1,Y).
\]
一个快速的「应用」如下：

\textbf{断言 12.2.} 若 $H\subseteq G$ 是子群，$X$ 是 $G$-模复形，则得到限制映射 $X^{hG}\to X^{hH}$（它在同伦意义下良定义）。

直观地说，$G$-不变的东西当然也是 $H$-不变的。

\textit{证明.} 考虑
\[
\Z[G/H]=\{f\colon G/H\to\Z\mid f\ \text{只在有限多个点处非零}\},
\]
它通过平移带有 $G$-作用，并且等同于从 $H$ 到 $G$ 由平凡模诱导出的模 $\Z[G]\otimes_{\Z[H]}\Z$。

\textbf{断言 12.3.} $\operatorname{Hom}^{\mathrm{der}}_G(\Z[G/H],X)\simeq X^{hH}$ 是拟同构。

\textit{证明.} 设 $P_H\xrightarrow{\mathrm{qis}}\Z$，其中 $P_H$ 是 $H$-模的投射复形。则得到诱导的 $G$-模 $\Z[G]\otimes_{\Z[H]}P_H$。注意 $\Z[G]$ 作为 $\Z[H]$-模是自由的：取 $G/H$ 的陪集代表元即得一基。因此 $\Z[G]\otimes_{\Z[H]}-$ 保持拟同构（我们已知 $\Z[H]$ 的情形，而 $\Z[G]$ 可看作 $\Z[H]$ 的若干拷贝的直和）。这意味着
\[
\Z[G]\otimes_{\Z[H]}P_H\xrightarrow{\mathrm{qis}}\Z[G/H],
\]
而后者作为 $\Z[G]$-模的复形是投射的：这是因为 $\Z[G]$ 与 $P_H$ 都有界，故它也有界，而诱导到 $\Z[G]$ 保持投射模，因为我们仍得到自由模的直和项。或者也可用下述万有性质：到无环复形的每个映射都零伦，而从诱导复形出发的 $G$-等变映射等同于从 $P_H$ 出发的 $H$-等变映射。于是由定义
\[
\operatorname{Hom}^{\mathrm{der}}_G(\Z[G/H],X):=\operatorname{Hom}_G(\Z[G]\otimes_{\Z[H]}P_H,X)=\operatorname{Hom}_H(P_H,X)=:X^{hH},
\]
即得断言。\hfill$\square$

由此可见，我们经映射
\[
\varepsilon\colon\Z[G/H]\to\Z,\qquad \sum_{g_i\in G/H}n_ig_i\mapsto\sum_i n_i
\]
得到映射 $X^{hG}\to X^{hH}$；当 $\Z$ 取平凡作用时，$\varepsilon$ 显然是 $G$-等变的。由前述讨论，我们得到导出函子之间的限制映射
\[
X^{hG}=\operatorname{Hom}^{\mathrm{der}}_G(\Z,X)\to\operatorname{Hom}^{\mathrm{der}}_G(\Z[G/H],X)=X^{hH},
\]
它在同伦意义下良定义（同伦意义下良定义，如此递推，即我们「惯常的错误」）。\hfill$\square$

回忆此处一切都是阿贝尔群的复形，故不存在「类型不匹配」。事实上，若 $H\le G$ 是有限指数的，则有 $G$-等变映射
\[
\Z\xrightarrow{\ \kappa\ }\Z[G/H]\xrightarrow{\ \varepsilon\ }\Z,\qquad 1\mapsto\sum_{g\in G/H}g,
\]
其复合对应乘以指数 $[G:H]$。这给出膨胀映射 $X^{hH}\to X^{hG}$，使得复合 $X^{hG}\to X^{hH}\to X^{hG}$ 同伦于乘以 $[G:H]$。

更具体地说，设有 $H$-不变的对象与 $G$-不变的对象。取 $G/H$ 的陪集代表元，便可对任一 $H$-不变元素取「相对范数」，得到一个 $G$-不变元素。这正是上面的映射所做的，也解释了复合为何乘以 $[G:H]$。

\textbf{定义 12.4.} $H^i(G,X):=H^i(X^{hG})$ 称为 $G$ 以 $X$ 为系数的（超）上同调。

前缀「超」过去用来指复形上的运算；若该复形只在次数 $0$ 有项，则称为「群上同调」。

\textbf{断言 12.5.} 若 $X$ 只在非负次数有项，即 $X^i=0$ 对一切 $i<0$，则 $H^0(G,X)=H^0(X)^G$，且 $H^i(G,X)=0$ 对 $i<0$。

\textit{证明.} 取 $\Z$ 作为 $G$-模的某个投射消解 $P$；由命题 11.9 的构造，可取 $P$ 只在非正次数有项。按定义，$H^0(G,X)$ 等同于映射 $f\colon P\to X$ 的同伦类，而所有这些映射都形如
\[
\begin{array}{ccccccccc}
\cdots & \longrightarrow & P^{-2} & \longrightarrow & P^{-1} & \longrightarrow & P^0 & \longrightarrow & 0\\
& & {\scriptstyle f}\big\downarrow & & \big\downarrow & & \big\downarrow & & \\
\cdots & \longrightarrow & 0 & \longrightarrow & 0 & \longrightarrow & X^0 & \longrightarrow & X^1
\end{array}
\]
于是 $f$ 的任何同伦都是 $0$，同理 $H^i(G,X)=0$ 对 $i<0$。由交换性必有 $df=fd=0$。由此可见，这样的映射 $f$ 等同于 $G$-等变映射
\[
\Z=P^0/dP^{-1}=\operatorname{Coker}(P^{-1}\to P^0)\to\ker(X^0\to X^1)=H^0(X),
\]
而后者按拟同构等同于 $H^0(X)$ 中的一个 $G$-不变向量（即经由 $1$ 的像给出）。\hfill$\square$

下面转向构造 $\Z$ 作为 $G$-模的显式投射消解。

\textbf{例 12.6.} 设 $G:=\Z/n\Z$，生成元为 $\sigma$。我们断言下述映射是拟同构：
\[
\begin{array}{ccccccccccc}
\cdots & \longrightarrow & \Z[G] & \xrightarrow{1-\sigma} & \Z[G] & \xrightarrow{\sum_i\sigma^i} & \Z[G] & \xrightarrow{1-\sigma} & \Z[G] & \xrightarrow{\varepsilon} & 0\\
& & \big\downarrow & & \big\downarrow & & \big\downarrow & & \big\downarrow & & \\
\cdots & \longrightarrow & 0 & \longrightarrow & 0 & \longrightarrow & 0 & \longrightarrow & \Z & \longrightarrow & 0
\end{array}
\]
负阶上同调的消没来自我们早先关于 Tate 上同调的结果；而 $\varepsilon$ 的核，即坐标和为零的元素的集合，正是 $1-\sigma$ 的像。

\textbf{推论 12.7.} 若 $M$ 是 $G$-模（视为次数 $0$ 处的复形），则 $M_{hG}$ 拟同构于复形
\[
\cdots\to0\to0\to M\xrightarrow{1-\sigma}M\xrightarrow{\sum_i\sigma^i}M\xrightarrow{1-\sigma}\cdots,
\]
其中第一个 $M$ 位于次数 $0$。

注意，从 $\Z[G]$ 到任一对象的 $G$-等变映射就是该对象的一个元素。确实，不变子模正是零阶上同调群，正如所愿。因此这一构造只给出了我们早先关于 Tate 上同调所学内容的「一半」。

现在我们想对每个 $G$ 构造显式消解。贯穿始终的「格言」是：「这类东西都来自 bar 构造。」设 $A$ 是交换环，$B$ 是 $A$-代数；最重要的情形是 $A:=\Z$、$B:=\Z[G]$。

\textbf{定义 12.8.} 对上述 $A$ 与 $B$，\textbf{bar 复形} $\operatorname{Bar}_A(B)$ 为
\[
\cdots\to B\otimes_AB\otimes_AB\xrightarrow{b_1\otimes b_2\otimes b_3\mapsto b_1b_2\otimes b_3-b_1\otimes b_2b_3}B\otimes_AB\xrightarrow{b_1\otimes b_2\mapsto b_1b_2}B\to0\to\cdots,
\]
其中 $B$ 位于次数 $0$。一般地，$\operatorname{Bar}^{-n}_A(B):=B^{\otimes_A n+1}$，微分为
\[
b_1\otimes\cdots\otimes b_{n+1}\mapsto b_1b_2\otimes b_3\otimes\cdots\otimes b_{n+1}
-b_1\otimes b_2b_3\otimes\cdots\otimes b_{n+1}
+b_1\otimes b_2\otimes b_3b_4\otimes\cdots\otimes b_{n+1}-\cdots.
\]

容易看出这一微分的平方为零：只需有选择地删去张量因子并检查符号，故它确实是链复形。

\textbf{断言 12.9.} $\operatorname{Bar}_A(B)$ 同伦等价于零。

\textit{证明.} 我们要找恒等映射的零伦，即映射 $h$ 使 $hd+dh=\mathrm{id}$：
\[
\begin{array}{ccccccccc}
\cdots & \longrightarrow & B\otimes_AB\otimes_AB & \longrightarrow & B\otimes_AB & \longrightarrow & B & \longrightarrow & 0\\
& & {\scriptstyle h_2}\big\uparrow{\scriptstyle\mathrm{id}} & & {\scriptstyle h_1}\big\uparrow{\scriptstyle\mathrm{id}} & & {\scriptstyle h_0}\big\uparrow{\scriptstyle\mathrm{id}} & & \\
\cdots & \longrightarrow & B\otimes_AB\otimes_AB & \longrightarrow & B\otimes_AB & \longrightarrow & B & \longrightarrow & 0
\end{array}
\]
于是定义 $h_0(b):=1\otimes b$，$h_1(b_1\otimes b_2):=1\otimes b_1\otimes b_2$。确实，此时有
\[
(dh_1+h_0d)(b_1\otimes b_2)=d(1\otimes b_1\otimes b_2)+1\otimes b_1b_2=b_1\otimes b_2-1\otimes b_1b_2+1\otimes b_1b_2=b_1\otimes b_2,
\]
正如所愿。容易证明，对一切 $n$ 类似地定义 $h_n$ 便给出恒等映射的零伦。\hfill$\square$

换个说法，考虑图
\[
\begin{array}{ccccccc}
\cdots & \longrightarrow & B\otimes_AB\otimes_AB & \longrightarrow & B\otimes_AB & \longrightarrow & 0\\
& & & & {\scriptstyle\mathrm{id}}\big\downarrow & & \big\downarrow{\scriptstyle b\otimes b'\mapsto bb'}\\
\cdots & \longrightarrow & 0 & \longrightarrow & B & \longrightarrow & 0
\end{array}
\]
其中微分中的 $d$ 就是乘法映射。这是同伦等价，因为它的锥正是 $\operatorname{Bar}_A(B)$。

把每一项看作双模（即左右各有交换作用的模），其中在首项上从左边乘以 $B$，在末项上从右边乘以 $B$。于是这些微分都是双模同态。给定（左）$B$-模 $M$，与 $M$ 在 $B$ 上作张量，便得到图
\[
\begin{array}{ccccccc}
\cdots & \longrightarrow & B\otimes_AB\otimes_AM & \longrightarrow & B\otimes_AM & \longrightarrow & 0\\
& & & & {\scriptstyle b\otimes m\mapsto bm}\big\downarrow & & \\
\cdots & \longrightarrow & 0 & \longrightarrow & M & \longrightarrow & 0
\end{array}
\]
它也是同伦等价（映射 $B\otimes_AM\to M$ 就是「作用映射」；并注意这些张量积有意义，因为 $B$ 是 $A$-模）。微分不变，只是末项换成了 $M$ 的元素，例如
\[
b_1\otimes b_2\otimes m\mapsto b_1b_2\otimes m-b_1\otimes b_2m.
\]
换言之，$M$ 典范地同伦等价于一个复形，其中每一项都形如 $B\otimes_AN$；此处 $N$ 代表 $B\otimes_A\cdots\otimes_AB\otimes_AM$，即由某个 $A$-模诱导出的模。

现在把这应用于 $A:=\Z$、$B:=\Z[G]$、$M:=\Z$（即平凡模）的情形。注意这只是 $B$-模复形之间的拟同构，而并非同伦等价：因为逆映射只是 $A$-线性的，而不是 $B$-线性的！确实，这样的逆映射应为
\[
B\otimes_AM\xrightarrow{b\otimes m\mapsto bm}M\xrightarrow{m\mapsto1\otimes m}B\otimes_AM,
\]
即 $b\otimes m\mapsto bm\mapsto1\otimes bm\neq b\otimes m=b(1\otimes m)$，因为作用在 $B$ 上而不在 $M$ 上，这并不是 $B$-模的作用。一般地，单向拟同构的存在并不蕴含反方向拟同构的存在；而同伦等价按定义包含反方向的映射，故是对称的。另外回忆，同伦等价蕴含拟同构。

于是我们得到 $\Z$ 的自由 $G$-模典范投射消解
\[
\begin{array}{ccccccccc}
\cdots & \longrightarrow & \Z[G^3] & \longrightarrow & \Z[G\times G] & \longrightarrow & \Z[G] & \longrightarrow & 0\\
& & \big\downarrow & & \big\downarrow & & \big\downarrow & & \\
\cdots & \longrightarrow & 0 & \longrightarrow & 0 & \longrightarrow & \Z & \longrightarrow & 0
\end{array}
\]
其微分为 $(g_1,g_2)\mapsto g_1g_2-g_1$，等等，因为 $G$ 平凡地作用于 $\Z$。注意 $\Z[G\times G]\cong\Z[G]\otimes_\Z\Z[G]$，因为两者都以乘积群的元素为基。

这是计算群上同调的一个很好的显式投射消解！最终得到形如
\[
\cdots\to0\to M\to\Z[G]\otimes M\to\Z[G\times G]\otimes M\to\cdots
\]
的复形，其中 $M$ 位于次数 $0$，且 $G$ 有限。
% [原文疑误：原印「位于 $\Z[G]\otimes M$ 的微分核中的元素」——按 $n$-上闭链的定义应为 $\Z[G^n]\otimes M$（该复形的第 $n$ 项是 $\Z[G^n]\otimes M$，其中 $\Z[G\times G]=\Z[G^2]$），此处径改]
位于 $\Z[G^n]\otimes M$ 的微分核中的元素称为以 $M$ 为系数的\textbf{群 $n$-上闭链}；位于微分像中的元素称为 \textbf{$n$-上边缘}。
